\documentclass[10pt]{article}

\usepackage[utf8]{inputenc}

\usepackage[T1]{fontenc}

\usepackage{amsmath}

\usepackage{amsfonts}

\usepackage{amssymb}

\usepackage[version=4]{mhchem}

\usepackage{stmaryrd}

\usepackage{mathrsfs}

\usepackage{graphicx}

\usepackage[export]{adjustbox}

\graphicspath{ {./images/} }

\usepackage{bbold}



\begin{document}

ТОПОЛОГИЧЕСКОЕ ПРОИЗВЕДЕНИЕ, $\mathrm{T}$ И $\mathrm{x}$ 0н о в с к о е п р о и з в е де в и е, семейства топологических пространств $\left\{\left(X_{\alpha}, \mathscr{\mathscr { V } \alpha}\right): \alpha \in A\right\}-$ топологич. пространство $(X, \mathscr{\sigma})$, где $X-$ декартово произведение (т. е. полное прямое произведение) множеств $X_{\alpha}$ по $\alpha \in A$ и $\mathscr{T}-$ слабейшая (т. е. наименьшая) топология на множестве $X$ такая, что все отображения естественного проектирования $\pi_{\alpha}:(X, \mathscr{\mathscr { T }}) \rightarrow\left(X_{\alpha}, \mathscr{\mathscr { \tau }} \alpha\right)$ непрерывны. Топология $\mathscr{T}$ наз. при этом т о п о л о г и е й



\includegraphics[max width=\textwidth, center]{2023_07_09_e2738812ec9849f41668g-1}

Іологич. произведением семейства пространств $\left\{\left(X_{\alpha}\right.\right.$, $\mathscr{\sigma} \alpha): \alpha \in A\}$



Стандартной базой топологич. пространства $(X, \mathscr{V})$ служит семейство всех множеств вида $\pi_{\alpha_{1}}^{-1}\left(U_{\alpha_{1}}\right) \cap \ldots \cap$ $\cap \pi_{\alpha_{n}}^{-1}\left(U_{\alpha_{n}}\right)$, где $\alpha_{1}, \ldots, \alpha_{n}-$ цроизвольный конечный набор элементов индексирующего множества $A$, a $U_{\alpha_{i}}$ - любой элемент топологии $\mathscr{F}_{\alpha_{i}}, \quad i=1, \ldots, n$.



В частности, если семейство $\left\{\left(X_{\alpha}, \mathscr{\sigma}_{\alpha}\right): \alpha \in A\right\}$ состоит из двух пространств $X$ и $Y$, то базу топологии произведения $Z=X \times Y$ образуют множества вида $U \times V$, где $U-$ произвольное открытое множество в $X$, a $V$ - произвольное открытое множество в $Y$. Аналогично описывается база топологии произведения любого конечного упорядоченного множества топологич. пространств. Примеры Т. п.: Ілоскость произведение двух прямых, $n$-мерное пространство $\mathbb{R}^{n}$ - произведение $n$ прямых, тор - произведение двух окружностей.



Первоначальные попытки определения Т. п. бесконечного множества топологич. пространств относились к случағо метризуемых сомножителей. Соответственно, топологию произведения пытались указать в терминах сходимости обычных (счетных) последовательностей. Когда семейство сомножителей несчетно, получить на этом пути тот же результат, что и выше, было невозможно, так как оператор замыкания в Т. п. несчетного множества неодноточечных метризуемых пространств не может быть описан на языке сходящихся последовательностей или сведен к замыканиям счетных множеств.



Определение Т. п. произвольного бесконечного множества топологич. пространств было дано А. Н. Тихоновым (1925). Он же доказал, что Т. п. бикомпактных пространств всегда является бикомпактным пространством.



Операция Т. п. является одним из главных средств формирования новых топологич. объектов из уже имею щихся. С ее помощью конструируется ряд основных стандартных объектов общей топологии - в частности, тихоновские кубы $I^{\tau}$, определяемые как топологич. произведение семейства мощности $\tau$ отрезков числовой прямой. По теореме Тихонова, все тихоновские кубы бикомпактны. А. Н. Тихонов доказал, что всякое вполне регулярное $T_{1}$-пространство гомеоморфно подпространству нек-рого куба $I^{\tau}$.



Кроме кубов $I^{\tau}$ важную роль в топологии играют пространства $D^{\tau}$ и $F^{\tau}$, являющиеся соответственно произведениями $\tau$ штук пространств, состоящих из двух изолированных точек (простых двоеточий $D$ ) и двухточечных пространств $\mathrm{c}$ одной изолированной точкой (связных двоеточий $F$ ). Всякий компакт есть непрерывный образ кавторова совершенного множества, т. е. пространства, гомеоморфного произведению счетного числа простых двоеточий $D$; всякое индуктивно нульмерное пространство, т. е. $T_{0}$-пространство, в к-ром открыто-замкнутые множества образуют базу, гомеоморфно подмножеству нек-рого канторова дисконтинуума $D^{\tau}$; всякое $T_{0}$-пространство гомеоморфно подмножеству пространства $F^{\tau}$. В связи с силой и ролью операции Т. п. центральное место в проблематике общей топологии занимают вопросы поведения топологич. свойств при операции Т. п. Устойчивы относительно операции Т. п. классы хаусдорфовых пространств, регулярных пространств и вполне регулярных пространств. Но произведение нормального пространства на отрезок может быть не нормальным пространством, не устойчивым относительно операции Т. п., даже в случае конечного числа сомножителей, такие важные топологич. свойства, как линделёфовость и паракомпактность.



Важную роль в общей топологии и ее приложениях (в частности, к построению моделей теории множеств) играет теорема: число Суслина Т. п. любого множества сепарабельных топологич. пространств счетно. В частности, счетно число Суслина любого тихоновского куба.



Операция Т. п. порождает, путем выделения во всем Т. п. определенных подпространств, весьма полезные операции $\Sigma$-произведения и $\sigma$-произведения. К совершенно иной топологии на декартовом произведении множеств $X_{\alpha}$ приводит операция ящичного произведения топологий $\mathscr{V} \alpha$.



Лum.: [1] А р х а н г е ль с к и й А. В., П о но м аp е в В. И., Основы общей топологии в задачах и упражнениях, М., 1974; [2] Б у р б а к и Н., Общая топология. Основные структуры, пер. с франц., М., 1968; [3] А л е к с а н д р о в П. С., Введение в теорию множеств и общую топологию, М., 1977.

$A . B$. Архангельский, B. В. Федориук.



ТОПОЛОГИЧЕСКОЕ ПРОСТРАНСТВО - совОКУПность двух объектов: множества $X$, состоящего из әлементов произвольной природы, наз. то чк а м и данного пространства, и из введенной в әто множество то п о ло г и ч е с к о й $\mathbf{~ т ~ р ~ у ~ к ~ т ~ у - ~}$ р ы, или т о ІІ о ло г и и, все равно - открытой или замкнутой (одна переходит в другую, если заменить множества, составляющие данную топологию, их дополнениями). Если не сказано противное, под топологией (\$) будет пониматься открытая топология. Логически самый простой способ определения топологии в данном множестве $X$ заключается в непосредственном указании тех подмножеств множества $X$, к-рые составляют эту топологию. Но часто проще определять не все множества, являющиеся элементами данной топологии, а только нек-рые множества этих әлементов (т. е. базу данной топологии), достаточные для того, чтобы все остальные әлементы топологии получались как объединения (в случав открытой топологии) или пересечения (в случае замкнутой) множеств, к-рые составляют базу. Так, напр., обычная топология числовой прямой получается, если определить в качестве базы ее открытой топологии множество всех интервалов (можно ограничиться даже одними интервалами с рациональными концами). Остальные открытые множества суть объединения интервалов.



Имея в виду базу открытой и, соответственно, замкнутой топологии, часто говорят об открытой, соответственно, замкнутой базе данного Т. П., причем открытые базы рассматриваются чаще замкнутых, поэтому если говорят просто о базе Т. п., то всегда имеют в виду его открытую базу. Наименьшее кардинальное число (в нетривиальных случаях бесконечное), являющееся мощностью какой-либо базы данного Т. п., наз. его весом. После мощности множества всех точек пространства вес является важнейшим т. н. к а р д и-



\includegraphics[max width=\textwidth, center]{2023_07_09_e2738812ec9849f41668g-1(1)}

с т р а н с т в а. Особенно важны простравства, имеющие счетную базу; напр., числовая прямая есть такое пространство. Аналогично, счетная база евклидова $n$-мерного пространства получается, если взять т.н. рациональные (открытые) шары, т. е. шары, радиус к-рых и координаты центра к-рых суть рациональные числа. Часто приходится определять тем или иным стандартным (естественным) способом топологию во





\end{document}